\section*{Étape B — Produit scalaire — Exercices type bac proposés}

\subsection*{Sujet 1 — Triangle, angle et orthogonalité}

Dans un repère orthonormé, on considère les points :
\[
A(1;2),\qquad B(5;4),\qquad C(3;-2).
\]
\begin{enumerate}
\item Calculer les coordonnées des vecteurs $\overrightarrow{AB}$ et $\overrightarrow{AC}$.
\item Calculer le produit scalaire $\overrightarrow{AB}\cdot\overrightarrow{AC}$.
\item Les droites $(AB)$ et $(AC)$ sont-elles perpendiculaires ?
\item Calculer les longueurs $AB$ et $AC$.
\item En déduire la valeur exacte de $\cos(\widehat{BAC})$.
\item Le triangle $ABC$ est-il rectangle ? Justifier.
\item Calculer $\overrightarrow{BA}\cdot\overrightarrow{BC}$ et interpréter le signe obtenu.
\end{enumerate}

\subsection*{Sujet 2 — Droite, vecteur normal et projeté orthogonal}

Dans un repère orthonormé, on considère la droite $d$ d'équation :
\[
2x-y-3=0
\]
et le point $A(4;1)$.
\begin{enumerate}
\item Donner un vecteur normal à la droite $d$.
\item Justifier que la droite $\Delta$, perpendiculaire à $d$ et passant par $A$, a pour vecteur directeur $\vec n(2;-1)$.
\item Déterminer une équation paramétrique de $\Delta$.
\item Déterminer les coordonnées du point $H$, intersection de $d$ et de $\Delta$.
\item Justifier que $H$ est le projeté orthogonal de $A$ sur $d$.
\item Calculer la distance $AH$.
\item Déterminer une équation cartésienne de la droite parallèle à $d$ passant par $A$.
\end{enumerate}

\subsection*{Sujet 3 — Cercle, diamètre et tangente}

Dans un repère orthonormé, on considère les points :
\[
A(-1;2),\qquad B(5;4).
\]
On note $\mathcal C$ le cercle de diamètre $[AB]$.
\begin{enumerate}
\item Déterminer les coordonnées du centre $I$ du cercle $\mathcal C$.
\item Calculer le rayon de $\mathcal C$.
\item Déterminer une équation cartésienne de $\mathcal C$.
\item Soit $M(x;y)$. Montrer que $M\in \mathcal C$ équivaut à :
\[
\overrightarrow{MA}\cdot\overrightarrow{MB}=0.
\]
\item Vérifier que le point $T(1;5)$ appartient au cercle $\mathcal C$.
\item Déterminer un vecteur normal à la tangente au cercle en $T$.
\item En déduire une équation de la tangente au cercle $\mathcal C$ en $T$.
\end{enumerate}
